\documentclass[uplatex,dvipdfmx,a4paper,11pt]{jsarticle}
\usepackage[top=25mm,bottom=25mm,left=23mm,right=23mm]{geometry}
\usepackage{longtable,booktabs,array}
\usepackage{enumitem}
\usepackage[dvipdfmx]{graphicx,xcolor}
\usepackage[dvipdfmx]{hyperref}
\usepackage{pxjahyper}
\usepackage{fancyhdr}
\usepackage{titlesec}
\setlist[itemize]{leftmargin=2em}
\setlength{\parindent}{1em}
\setlength{\parskip}{0.4em}
\pagestyle{fancy}
\fancyhf{}
\lhead{第10章 ソフトウェア工学プロセス}
\rhead{SWEBOK v4.0a 日本語まとめ}
\cfoot{\thepage}
\newcommand{\pointbox}[2]{\begin{quote}\fbox{\begin{minipage}{0.92\linewidth}\textbf{#1}\par #2\end{minipage}}\end{quote}}
\newcommand{\todobox}[2]{\begin{quote}\fbox{\begin{minipage}{0.92\linewidth}\textbf{TODO（図）} #1\par\smallskip\textbf{画像生成用プロンプト}\par #2\end{minipage}}\end{quote}}
\title{第10章 ソフトウェア工学プロセス\\\large Software Engineering Process の日本語まとめ\\\normalsize SWEBOK Guide v4.0a Chapter 10 を初学者向けに再構成}
\date{作成日：2026年5月12日}
\begin{document}
\maketitle
\thispagestyle{empty}
\section*{この資料の主張}
ソフトウェア工学プロセスは、開発手順の名前を選ぶだけの話ではない。入力を成果物へ変換する活動、活動を束ねるライフサイクル、経験的測定に基づく監視と改善を、対象システムと組織の文脈に合わせて設計・運用することが中核である。
\section*{この資料の読み方}
専門用語は原則として日本語で示し、必要に応じて英語を括弧内に併記する。初学者が前提知識なしに読めるように、各節では「何のための概念か」「実務では何を見ればよいか」を重視する。
\section*{出典と位置づけ}
原資料：IEEE Computer Society, Guide to the Software Engineering Body of Knowledge (SWEBOK Guide) v4.0a, Chapter 10, Software Engineering Process. 本資料は学習用の日本語要約であり、公式翻訳ではない。
\tableofcontents
\clearpage
\section{全体概要：この章で学ぶこと}
この章は、ソフトウェアを作り、運用し、保守し、最後に廃止するまでの「作業の組み立て方」を扱う。ここでいうプロセスは、プログラムが実行中に進む処理の流れではなく、人やチームが成果物を作るために行う作業活動である。

この章の中心は三つである。第一に、プロセスとは何かという基礎概念である。第二に、複数のプロセスを時間軸や反復構造の中に配置するライフサイクルである。第三に、実行中のプロセスを測定し、評価し、改善する方法である。

初学者は、ウォーターフォール、アジャイル、DevOps といった名前を暗記する前に、「入力を活動で変換して出力を作る」「出力は別の活動の入力になる」「すべてのプロジェクトに万能なプロセスはない」「測定なしにプロセスの良し悪しは判断しにくい」という四点を押さえるとよい。

\pointbox{到達目標}{この章を読み終えると、プロセス、活動、タスク、ライフサイクル、ライフサイクルモデル、プロセス評価、プロセス改善の違いを説明できる。また、予測型、反復型、進化型、増分型、継続的開発、アジャイル、DevOps を、単なる流行語ではなく、要求の変化、納品頻度、測定、組織連携という観点から比較できるようになる。}
\todobox{10章の全体地図を入れる。}{白背景、モノクロ、A4本文幅に収まる横長の学習用図解を作成する。中央に「ソフトウェア工学プロセス」を置き、周囲に「基礎」「ライフサイクル」「プロセス評価と改善」を大きな三領域として配置する。基礎には「入力・活動・出力」、ライフサイクルには「概念から廃止まで」、評価と改善には「測定・PDCA・GQM」を入れる。日本語ラベル、細線、講義ノート風、余白を広めにする。}
\section{最初に必要な用語}
\paragraph{表1：この章で最初に押さえる用語}
\begin{longtable}{>{\raggedright\arraybackslash}p{0.28\linewidth}>{\raggedright\arraybackslash}p{0.64\linewidth}}
\toprule
用語 & 初学者向けの説明\\
\midrule
\endfirsthead
\toprule
用語 & 初学者向けの説明\\
\midrule
\endhead
プロセス & 入力を受け取り、作業活動によって出力へ変換するまとまりである。ここでは人間の作業活動を指す。\\
活動 & プロセスを構成するまとまった作業である。例として設計、実装、検証がある。\\
タスク & 活動をさらに細かくした、実行すべき作業単位である。\\
ライフサイクル & 構想から廃止まで、システムや製品がたどる全体の変化である。\\
ライフサイクルモデル & 作業をどの順序で、どの程度反復し、どの時点で成果物を出すかを示す型である。\\
プロジェクト & 明確な開始と終了を持ち、固有の製品・サービス・結果を作る一時的な取り組みである。\\
測定 & 状態や成果を数値や指標で把握する活動である。プロセス改善の土台になる。\\
経験的証拠 & 実際の測定や観察から得た根拠である。プロセス選択や改善判断を支える。\\
\bottomrule
\end{longtable}
\section{導入：プロセスを学ぶ理由}
ソフトウェア工学プロセスは、開発手順の名前を選ぶだけの問題ではない。要求を理解し、設計し、実装し、検証し、移行し、運用し、保守し、廃止するまでには、多数の活動が相互に関係する。ある活動の出力は、別の活動の入力になる。したがって、単独の作業だけでなく、作業同士の接続を設計する必要がある。

SWEBOK 第10章は、プロセスを「概念」「ライフサイクル」「プロセス評価と改善」の三つの視点から整理する。ソフトウェア工学の世界では、プロセス標準化に関する議論が活発であり、ISO/IEC/IEEE 12207 のような標準も重要な参照点になる。

重要なのは、理想的な唯一のプロセスは存在しないという点である。プロセスは、プロジェクトの性質、製品の規模、利害関係者、組織の成熟度、変更の頻度、求められる品質水準に合わせて選択・適応・運用される。さらに、その運用は経験的測定によって支えられるべきである。

\pointbox{この章の読み方}{本資料では、専門用語は原則として日本語で示し、必要に応じて英語を括弧内に併記する。英語名は、元資料やツール上の名称を確認するための手がかりである。説明は初学者向けに再構成しているため、原文の順序を保ちつつ、実務での見方を補足している。}
\section{略語}
\paragraph{表2：第10章に出てくる主な略語}
\begin{longtable}{>{\raggedright\arraybackslash}p{0.16\linewidth}>{\raggedright\arraybackslash}p{0.30\linewidth}>{\raggedright\arraybackslash}p{0.46\linewidth}}
\toprule
略語 & 日本語 & 説明\\
\midrule
\endfirsthead
\toprule
略語 & 日本語 & 説明\\
\midrule
\endhead
BPMN & 業務プロセスモデリング表記法 & 業務や開発作業の流れを図で表すための記法である。\\
CASE & コンピュータ支援ソフトウェア工学 & 開発活動をツールで支援する考え方や環境である。\\
CMM & 能力成熟度モデル & 組織のプロセス能力や成熟度を段階的に評価するモデルである。\\
CMMI & 能力成熟度モデル統合 & CMM 系の考え方を統合したプロセス改善モデルである。\\
GQM & 目標・質問・メトリクス & 測定を、目標から質問、質問から測定値へ落とし込む方法である。\\
IDEF0 & 統合定義 0 & 機能や活動の入出力、制約、支援資源を表すプロセス記法である。\\
PDCA & 計画・実行・確認・改善 & プロセスを経験的に改善するための基本循環である。\\
SLCM & ソフトウェアライフサイクルモデル & 作業の並べ方や反復の仕方を示す型である。\\
SLCP & ソフトウェアライフサイクルプロセス & ライフサイクル中で実行される個々のプロセスである。\\
UML & 統一モデリング言語 & 構造や振る舞いをモデルとして表す汎用的な記法である。\\
\bottomrule
\end{longtable}
\section{1 ソフトウェア工学プロセスの基礎}
\subsection{1.1 導入}
ソフトウェア工学プロセスとは、ソフトウェアを構築し、運用するためにソフトウェア技術者が実施する作業活動である。一般的な工学プロセスは、資源を消費しながら、入力を出力へ変換する相互に関連した活動の集合である。ソフトウェア工学も工学の一分野であるため、設計、製造、運用といった他の工学分野の考え方から多くを学んできた。

ただし、ソフトウェアでは、化学反応器のように同一の物理製品を多数製造するわけではない。むしろ、多数のソフトウェア単位を構築し、それらの流れを管理する点で製造の比喩が役立つ。さらに、ソフトウェア単位の実行は、ある種類のデータを別の種類のデータへ変換する操作として理解できる。

この章で使う「プロセス」という語は、プログラムが実行時に動く処理ではなく、技術者やチームが行う作業活動を意味する。プロセスは多くの知識領域と関係するが、とくにソフトウェア工学マネジメント、モデルと手法、品質、アーキテクチャ、テスト、測定、根本原因分析と深く関係する。

\pointbox{重要}{どのプロジェクトにもそのまま適用できる理想プロセスは存在しない。プロセスは、プロジェクトと組織の文脈に合わせて選択し、適応し、測定しながら運用するものである。}
\todobox{プロセスの入力・活動・出力の図を入れる。}{白背景、モノクロ、A4本文幅の横長図を作成する。左に「入力」、中央に「活動」、右に「出力」を置き、矢印で接続する。活動の上に「制約・方針」、下に「人員・ツール・インフラ」を配置し、活動を支えることを示す。出力から次のプロセスの入力へ戻る点線矢印も入れる。日本語ラベル、細線、講義ノート風。}
\subsection{1.2 ソフトウェア工学プロセスの定義}
プロセスは、相互に関係し合う活動の集合であり、入力を出力へ変換する。活動は、まとまりのあるタスクの集合である。タスクは、一つ以上のプロセス成果を達成するために必要、推奨、または許容される行動である。

プロセスの説明には、通常、必要な入力、変換活動、生成される出力が含まれる。さらに、指示、規則、制約のような制御要素と、ツール、技術、人員、インフラのような支援要素も関係する。あるプロセスの出力は、別のプロセスの入力になることが多い。

ソフトウェアシステムは、ソフトウェアだけでなく、ハードウェア、人、手作業の手順を含むことがある。したがって、ソフトウェア工学プロセスを考えるときは、ソフトウェア部分だけでなく、その周辺の人と組織の動きも含めて捉える必要がある。

\paragraph{表3：プロセスを記述するときに見る要素}
\begin{longtable}{>{\raggedright\arraybackslash}p{0.16\linewidth}>{\raggedright\arraybackslash}p{0.30\linewidth}>{\raggedright\arraybackslash}p{0.46\linewidth}}
\toprule
要素 & 意味 & 例\\
\midrule
\endfirsthead
\toprule
要素 & 意味 & 例\\
\midrule
\endhead
入力 & 活動の前提になる成果物や情報である。 & 要求、標準、既存コード、設計方針。\\
活動 & 入力を変換する作業である。 & 分析、設計、実装、レビュー、テスト。\\
出力 & 活動の結果として作られる成果物である。 & 設計書、コード、テスト結果、運用手順。\\
制御 & 活動を縛る方針や条件である。 & 契約、規制、品質基準、納期制約。\\
支援機構 & 活動を可能にする資源である。 & 人員、ツール、開発環境、組織知識。\\
\bottomrule
\end{longtable}
\section{2 ライフサイクル}
\subsection{2.1 ライフサイクルの定義、プロセス分類、用語}
ライフサイクルとは、システム、製品、サービス、プロジェクトなどが、構想から廃止まで進化していく全体である。ソフトウェア工学では、単に必要なプロセスを列挙するだけでは複雑さを十分に説明できないため、複数のプロセスと制約を含むライフサイクルとして捉える。

開発とは、利害関係者のニーズに従ってソフトウェアシステムを構築または変更する重要な段階である。製品ライフサイクルは、構想、提供、成長、成熟、廃止へ進む製品の変化として説明できる。これはソフトウェア固有ではなく、製品一般に当てはまる考え方である。

ソフトウェアシステムには、単体テストの対象になり得る原子的なソフトウェア構成要素、すなわちソフトウェア単位が含まれる。システム全体のライフサイクルは、発想、取得、供給、開発、運用、進化、保守、廃止までを含む。

\paragraph{表4：ISO/IEC/IEEE 12207 に基づくプロセス分類の学習用整理}
\begin{longtable}{>{\raggedright\arraybackslash}p{0.16\linewidth}>{\raggedright\arraybackslash}p{0.30\linewidth}>{\raggedright\arraybackslash}p{0.46\linewidth}}
\toprule
分類 & 役割 & 主な例\\
\midrule
\endfirsthead
\toprule
分類 & 役割 & 主な例\\
\midrule
\endhead
技術プロセス & ソフトウェア製品を作り、進化させ、運用し、廃止するための工学的活動である。 & 業務分析、利害関係者ニーズ定義、要求定義、アーキテクチャ、設計、実装、統合、検証、妥当性確認、移行、運用、保守、廃止。\\
技術管理プロセス & プロジェクトを計画し、制御し、リスクや品質を管理する活動である。 & プロジェクト計画、評価と制御、意思決定、リスク管理、構成管理、情報管理、測定、品質保証。\\
組織的プロジェクト支援プロセス & 組織としてプロジェクトを支える基盤を整える活動である。 & ライフサイクルモデル管理、インフラ管理、ポートフォリオ管理、人材管理、品質管理、知識管理。\\
合意プロセス & 取得・供給など、組織間または関係者間の合意を支える活動である。 & 取得プロセス、供給プロセス。\\
\bottomrule
\end{longtable}
\todobox{四つのプロセスカテゴリの図を入れる。}{白背景、モノクロ、A4本文幅の概念図を作成する。中央に「ソフトウェアシステムを生み出すライフサイクル」を置き、周囲に「技術プロセス」「技術管理プロセス」「組織的プロジェクト支援プロセス」「合意プロセス」の四つの箱を配置する。各箱に代表例を二、三個入れる。四つの箱から中央へ矢印を向ける。日本語ラベル、講義ノート風。}
\subsection{2.2 ライフサイクルが必要な理由}
ソフトウェア製品の作成、運用、廃止には、多数のプロセス、活動、タスク、制約が関わる。さらに、ソフトウェアシステムは、人、手順、ハードウェアを含むため、単純な工程表だけでは表現しにくい。

プロセスの入力は、多くの場合、別のプロセスの出力である。したがって、各プロセスは独立しているのではなく、相互に依存している。この相互依存が、ソフトウェア工学プロセス全体を複雑にする。

ライフサイクル仕様は、工学的な取り組みを可能にする強力な道具である。ライフサイクルを定義することは、各プロセスと制約を定義することであり、関係者間のコミュニケーション、技術管理、測定、評価、改善、品質管理の土台になる。

\pointbox{実務での見方}{「どの工程を先にやるか」だけでなく、「どの活動の出力を誰が使うか」「何を測れば状態を把握できるか」「変更が起きたときにどこへ戻るか」を考えると、ライフサイクルを実務的に理解しやすい。}
\subsection{2.3 プロセスモデルとライフサイクルモデルの概念}
プロセスモデルは、活動、入力、出力、制約、関係をどのように表すかというモデルである。これに対して、ライフサイクルモデルは、活動をどの順序や反復構造で配置するかを示す型である。

プロジェクト固有の活動列は、標準で定義された活動を、選択したソフトウェアライフサイクルモデルに写像することで作られる。つまり、現実のプロジェクトのライフサイクルは、標準や組織規程をそのまま丸写しするものではなく、選んだモデルに従って具体化するものである。

代表的なライフサイクルモデルには、ウォーターフォールモデル、Vモデル、増分モデル、スパイラルモデル、アジャイルモデルがある。

\paragraph{表5：プロセスモデルとライフサイクルモデルの違い}
\begin{longtable}{>{\raggedright\arraybackslash}p{0.16\linewidth}>{\raggedright\arraybackslash}p{0.30\linewidth}>{\raggedright\arraybackslash}p{0.46\linewidth}}
\toprule
観点 & プロセスモデル & ライフサイクルモデル\\
\midrule
\endfirsthead
\toprule
観点 & プロセスモデル & ライフサイクルモデル\\
\midrule
\endhead
主な関心 & 活動が何を入力とし、何を出力するか。 & 活動をどの順序、段階、反復で実施するか。\\
例 & IDEF0、BPMN、UML活動図、データフロー図。 & ウォーターフォール、Vモデル、増分、スパイラル、アジャイル。\\
実務での使い方 & 作業の中身と接続を説明する。 & プロジェクト全体の進み方を説明する。\\
\bottomrule
\end{longtable}
\subsection{2.4 開発ライフサイクルモデルの考え方}
各ソフトウェアシステムには、利害関係者の業務的・技術的ニーズを反映した固有の特徴がある。したがって、適切なライフサイクルは、その特徴を考慮して決める必要がある。

予測型ライフサイクルでは、範囲、時間、費用を早い段階で決める。これは、実装される要求集合が実質的に閉じており、大きな変更が起きないことを前提にする。

反復型ライフサイクルでは、範囲は比較的早く定まるが、製品理解が深まるにつれて時間や費用の見積りを更新する。進化型ライフサイクルでは、要求や顧客ニーズの変化、または要求の段階的導入により、製品やサービスが寿命を通じて変わる。増分型ライフサイクルでは、決められた期間の反復を通じて成果物へ機能を順次追加する。

継続的開発は、自動化ツールを用いて、新しいシステムやソフトウェアをステージング環境や各種テスト環境へ頻繁にリリースできるようにする実践群である。要求仕様の変更を許すライフサイクルは、変更に開かれているといえる。

\paragraph{表6：開発ライフサイクルの考え方の比較}
\begin{longtable}{>{\raggedright\arraybackslash}p{0.16\linewidth}>{\raggedright\arraybackslash}p{0.30\linewidth}>{\raggedright\arraybackslash}p{0.46\linewidth}}
\toprule
考え方 & 特徴 & 初学者向けの見分け方\\
\midrule
\endfirsthead
\toprule
考え方 & 特徴 & 初学者向けの見分け方\\
\midrule
\endhead
予測型 & 範囲・時間・費用を早期に決め、要求の大幅変更を想定しにくい。 & 最初に決めた計画への適合を重視する。\\
反復型 & 同じ種類の活動を繰り返し、理解の進展に応じて見積りを更新する。 & 作りながら学び、次の反復へ反映する。\\
進化型 & 製品やサービスが寿命を通じて変化する。 & ニーズや技術の変化に合わせて育てる。\\
増分型 & 反復ごとに機能を追加し、最後に十分な能力を持つ成果物になる。 & 小さな成果を積み上げる。\\
継続的開発 & 自動化により、リリースやテスト環境への反映を高頻度に行う。 & 短い間隔で出し、測り、直す。\\
\bottomrule
\end{longtable}
\todobox{ライフサイクルモデル比較図を入れる。}{白背景、モノクロ、A4本文幅の比較表風図を作成する。列に「予測型」「反復型」「進化型」「増分型」「継続的開発」を置く。行に「要求変更への姿勢」「成果物の出し方」「見積りの扱い」「向く状況」を置く。各セルは短い日本語。講義ノート風、細線、印刷で読みやすい。}
\subsection{2.5 開発ライフサイクルモデルと工学的側面}
ウォーターフォールモデルは、ソフトウェア工学の初期に広く知られるようになった予測型のモデルである。要求、予備設計、詳細設計、コーディング、テストといったフェーズを順序立てて進め、前のフェーズが完了するまで次に進まない厳格な進め方を重視する。文書駆動型であり、ソフトウェア危機に対する初期の工学的対応として意味を持った。

Vモデルは、ウォーターフォール型の変種または拡張として理解できる。増分型ライフサイクルでは、工程が完全な直列ではなく重なり合いながら進む。増分型であっても、要求を先に閉じるなら予測型にもなり得る。

スパイラルモデルは、進化型でリスク駆動のモデルである。要求分析、設計、コーディング、統合、テストなどを反復し、リスクを見ながら完成へ近づける。迅速プロトタイピングは、早期に試作品を作り、早いフィードバックを得ることを重視する。統一プロセス、RUP、OpenUP は、反復型・増分型のプロセス枠組みとして知られる。

アジャイル宣言は、ソフトウェア工学コミュニティの考え方を大きく変えた。アジャイルは、要求が利用者ニーズの変化に応じてどの段階でも変更され得ること、チームと顧客の信頼、コミュニケーション、価値の継続的提供、技術的卓越性を重視する。重要なのは、アジャイルは文書を不要とするものではないという点である。必要な文書は必要である。

アジャイルに関する誤解も多い。アジャイルは単一の方法ではない。文書を作らないから速い、という意味でもない。エクストリームプログラミング、スクラム、スプリント計画、ふりかえり、ペアプログラミングなど、多様な方法と実践が存在する。大規模プロジェクトやポートフォリオへアジャイルを拡張することは、いまなお難しい課題である。

DevOps は、開発、運用、関連する利害関係者のコミュニケーションと協力を高め、ソフトウェアとシステムの構築、パッケージ化、デプロイ、運用、改善をつなげる考え方である。より頻繁なリリースを可能にすることは、適切なプロセス管理のもとでは競争力になる。

ウォーターフォール対アジャイルのような論争は、歴史的背景を持つ。ソフトウェア工学では、好みや流行ではなく、経験的測定に基づいて議論すべきである。プロセスや製品の測定は、意思決定に欠かせない。

\todobox{代表的ライフサイクルモデルの位置づけ図を入れる。}{白背景、モノクロ、A4本文幅の二軸図を作成する。横軸を「要求変更に閉じている → 変更に開いている」、縦軸を「単一納品に近い → 頻繁なリリースに近い」とする。点として「ウォーターフォール」「Vモデル」「増分」「スパイラル」「迅速プロトタイピング」「統一プロセス」「アジャイル」「DevOps」を配置する。日本語ラベル、講義ノート風。}
\subsection{2.6 SLCPの管理}
SLCP はソフトウェアライフサイクルプロセスである。一般的なライフサイクルには、概念、開発、生産、利用、支援、廃止という段階がある。

概念段階では、利害関係者のニーズを識別し、概念を探索し、解決策を提案する。開発段階では、利用者ニーズを表す要求を洗練し、解決策を作り、システムを構築し、必要な検証と妥当性確認を行う。生産段階は対象システムの性質により範囲が異なるが、一般にはシステムの生産とテストを含む。利用段階では、システムが利用者ニーズを満たすために運用される。支援段階では、満足できる運用を実現するための支援行為を行う。廃止段階では、定められた手順に従ってシステムを処分する。

これらの段階は、必ず直列に進むわけではない。実際のライフサイクル仕様では、段階間の遷移を定義する。個別のプロジェクトでは、利害関係者にとって意味のある固有の段階が設定され、それらが一般的段階に対応づけられる。

\paragraph{表7：一般的なライフサイクル段階}
\begin{longtable}{>{\raggedright\arraybackslash}p{0.28\linewidth}>{\raggedright\arraybackslash}p{0.64\linewidth}}
\toprule
段階 & 主な内容\\
\midrule
\endfirsthead
\toprule
段階 & 主な内容\\
\midrule
\endhead
概念 & ニーズを識別し、概念と解決策を探索する。\\
開発 & 要求を洗練し、解決策を作り、検証と妥当性確認を行う。\\
生産 & 対象に応じて、システムの生産とテストを行う。\\
利用 & システムを運用し、利用者ニーズを満たす。\\
支援 & 満足できる運用のための保守・支援行為を行う。\\
廃止 & 定められた手順でシステムを処分する。\\
\bottomrule
\end{longtable}
\todobox{ライフサイクル段階と遷移の図を入れる。}{白背景、モノクロ、A4本文幅の横長図を作成する。「概念」「開発」「生産」「利用」「支援」「廃止」の箱を並べる。単純な直線矢印だけでなく、利用から支援、支援から開発へ戻る点線矢印も入れ、段階が必ず一方向ではないことを示す。日本語ラベル、講義ノート風。}
\subsection{2.7 ソフトウェア工学プロセス管理}
プロセス管理とは、製品を開発したりサービスを実施したりする作業を、方向づけ、制御し、調整することである。ソフトウェア工学プロセスは複数の管理階層で統制される。

最下層には、要求定義、設計、実装、検証、運用、保守といった技術プロセスがある。第二層には、プロジェクト計画やリスク管理などを含む技術管理レベルがある。第三層には、知識管理、ライフサイクルモデル管理、ポートフォリオ管理などを扱う経営または組織的支援レベルがある。

この階層を区別すると、問題の所在を見誤りにくくなる。実装が遅れている原因が、技術者の作業ではなく、プロジェクト計画、ツール統合、人材配置、組織知識の不足にあることもある。

\subsection{2.8 ソフトウェアライフサイクル適応}
各ソフトウェアシステムには固有の特徴がある。製品規模、複雑さ、利害関係者、規制、品質水準、運用環境、組織文化が異なれば、適切なライフサイクルも異なる。

ライフサイクル適応では、関連する特徴を識別し、適切な標準や組織内文書を選び、開発戦略やライフサイクルモデル、段階、プロセスを選択する。そして、その判断と根拠を文書化する。

ここで重要なのは、標準に書かれた名前をすべて保つ必要はなく、すべてのプロセスをそのまま含める必要もないという点である。標準は土台であり、現実のプロジェクトでは適用範囲と調整理由を明確にすることが重要である。小規模組織向けの ISO/IEC 29110 系列は、12207 から派生した例として理解できる。

\pointbox{適応の勘所}{「標準に従う」とは、標準を機械的に写すことではない。何を採用し、何を省き、なぜそうしたかを説明できる状態にすることである。}
\subsection{2.9 実務上の考慮}
ライフサイクルプロセスを定義するとは、技術プロセス、組織的プロセス、技術管理プロセス、合意プロセスを具体化することである。つまり、何を作るかだけでなく、誰が支え、どう管理し、どのように合意するかを定義する。

ソフトウェア工学は誕生以来、学び続けてきた分野である。製品の複雑さは増し続けているため、製品特性と外部特性の両方を考慮する必要がある。利害関係者の特徴も、ライフサイクル選択に影響する。

見積りと測定は不可欠である。ライフサイクルの中で誤った見積りや不確実な見積りを使うと、プロジェクト失敗につながる。ただし、正確な見積りは容易ではない。したがって、測定値には精度と不確実性を明示し、プロセスと製品の状態を正確に把握する必要がある。

現在の傾向として、継続的デリバリーが重視される。大きなプロセスを長期間進めながら途中成果物を出さないと、不確実性が増える。アジャイルの考え方は、途中で価値を示し、コミュニケーションを高め、不確実性を下げる点に貢献している。

\pointbox{論争より測定}{予測型かアジャイルかという議論は、現実のプロセス測定と製品測定に基づけるべきである。測定が不確実なら、プロセス選択そのものを振り返る必要がある。}
\subsection{2.10 ソフトウェアプロセス基盤、ツール、手法}
ソフトウェアプロセスを定義する記法には、自然言語、活動とタスクの一覧、データフロー図、状態図、IDEF0、ペトリネット、UML活動図、BPMN などがある。記法を選ぶ目的は、プロセスを人が理解し、合意し、管理し、改善できるようにすることである。

プロセス基盤には、プロセス定義を支援するツールだけでなく、テスト、構成管理、チケット管理など、個別プロセスを支援するツールが含まれる。重要なのは、ツールが単独で存在するのではなく、互いに統合されることである。

たとえば、あるコード単位がテストに合格したなら、その情報はテストツールだけに閉じていては不十分である。構成管理ツールやチーム全体の状態把握に連携されることで、次の活動が正しく判断できる。

統合されたツール群は、ソフトウェア工学環境と呼ばれることがある。1980年代から1990年代には CASE という語がよく使われたが、万能薬として過度に期待された面もあった。現在でも、バージョン管理、テスト、チケット管理などの自動化は、成功するライフサイクル実装に不可欠である。

\paragraph{表8：プロセス定義に使われる表現例}
\begin{longtable}{>{\raggedright\arraybackslash}p{0.28\linewidth}>{\raggedright\arraybackslash}p{0.64\linewidth}}
\toprule
表現 & 向いている用途\\
\midrule
\endfirsthead
\toprule
表現 & 向いている用途\\
\midrule
\endhead
自然言語 & 関係者が読みやすい説明。曖昧さには注意が必要である。\\
活動・タスク一覧 & 作業を漏れなく列挙したい場合。\\
データフロー図 & 情報や成果物の流れを示したい場合。\\
状態図 & 状態と遷移を明確にしたい場合。\\
IDEF0 & 入力、制約、出力、支援資源を整理したい場合。\\
UML活動図 & 活動の分岐、合流、並行を示したい場合。\\
BPMN & 業務プロセスや組織横断の流れを示したい場合。\\
\bottomrule
\end{longtable}
\todobox{ツール連携としてのソフトウェア工学環境の図を入れる。}{白背景、モノクロ、A4本文幅のネットワーク図を作成する。中央に「ソフトウェア工学環境」を置き、周囲に「バージョン管理」「テスト」「構成管理」「チケット管理」「ビルド」「要求管理」「測定ダッシュボード」を配置する。各ツール間を細線でつなぎ、「情報を共有する」ことを示す。日本語ラベル、講義ノート風。}
\subsection{2.11 プロセス監視とソフトウェア製品の関係}
開発者や管理者は、ソフトウェア工学プロセスの実行を監視し、プロセス目標が満たされているか、リスクが高まっていないかを評価する必要がある。これはプロセス評価の一部であり、ソフトウェア工学マネジメントとも関係する。

経験的方法は、プロセスの評価と改善だけでなく、製品の評価と改善も支える。プロセス実行の目的は、利害関係者のニーズを満たす製品を得ることである。したがって、プロセスだけを測っても不十分であり、製品の状態とあわせて見る必要がある。

たとえば、計画通りにレビュー回数が実施されていても、欠陥が減っていないなら、プロセスは目的を果たしていない可能性がある。逆に、製品品質が高く見えても、チームが過負荷で属人的な作業に依存しているなら、長期的なプロセスリスクがある。

\todobox{プロセス監視と製品評価の関係図を入れる。}{白背景、モノクロ、A4本文幅の循環図を作成する。「プロセス実行」から「成果物・製品」へ矢印を引く。プロセス側に「進捗、欠陥検出率、リードタイム」、製品側に「品質、欠陥、利用者価値」を配置する。両方から「測定・分析」へ矢印を集め、そこから「改善判断」へ、さらにプロセス実行へ戻す。日本語ラベル、講義ノート風。}
\section{3 ソフトウェアプロセスの評価と改善}
\subsection{3.1 ソフトウェアプロセス評価と改善の概要}
実行されたプロセスは改善できる、という考え方は、シェwhart＝デミングの PDCA、すなわち計画・実行・確認・改善の循環に表れている。ソフトウェア工学プロセスでも、プロセスを評価し、改善するための多くの方法が発展してきた。

PDCA の重要性は、意思決定に経験的証拠を使う点にある。どのライフサイクルを選ぶか、プロセスをどう評価するか、どの改善を行うかといった判断は、観察や測定に基づくべきである。

プロセス評価は、監査のためだけの活動ではない。目的は、チームや組織がよりよい成果物を安定して作れるようにすることである。評価結果が改善行動につながらなければ、測定は形式的な作業になってしまう。

\todobox{PDCA によるプロセス改善の図を入れる。}{白背景、モノクロ、A4本文幅の循環図を作成する。「計画」「実行」「確認」「改善」の四つの箱を円形に配置する。中央に「経験的証拠に基づく意思決定」と書く。各段階に短い説明を添える。日本語ラベル、講義ノート風。}
\subsection{3.2 GQM：目標・質問・メトリクス}
GQM は Goal-Question-Metric の略であり、日本語では目標・質問・メトリクスと呼べる。測定を闇雲に始めるのではなく、まず何を達成したいかという目標を定める。次に、その目標が達成されたかを判断するための質問を作る。最後に、その質問に答えるための測定値を定める。

GQM は Basili の品質改善パラダイムに基づく。測定可能な目標を設定し、何かを変え、その変化の効果を評価する。評価が肯定的であれば、改善が起きたと判断できる。

初学者が誤解しやすい点は、メトリクスが先ではないということだ。測れるものを集めるのではなく、改善したい目標から逆算して測るべきものを決める。

\paragraph{表9：GQM の例}
\begin{longtable}{>{\raggedright\arraybackslash}p{0.28\linewidth}>{\raggedright\arraybackslash}p{0.64\linewidth}}
\toprule
段階 & 例\\
\midrule
\endfirsthead
\toprule
段階 & 例\\
\midrule
\endhead
目標 & リリース前に重大欠陥をより早く見つけたい。\\
質問 & 重大欠陥はどの工程で検出されているか。レビューで検出できる割合は増えているか。\\
メトリクス & 重大欠陥数、工程別検出件数、レビュー検出率、修正リードタイム。\\
\bottomrule
\end{longtable}
\todobox{GQM の階層図を入れる。}{白背景、モノクロ、A4本文幅の縦長図を作成する。最上段に「目標」、中段に複数の「質問」、下段に複数の「メトリクス」を配置し、上から下へ枝分かれする矢印で接続する。右側に「測定は目標から逆算する」と注記する。日本語ラベル、講義ノート風。}
\subsection{3.3 枠組みに基づく方法}
プロセス評価方法の中には、プロセス参照モデルと評価参照モデルを使う枠組みに基づくものがある。代表例には、CMM、CMMI、ISO/IEC 33000 系列、いわゆる SPICE がある。

ISO/IEC 33000 枠組みは、プロセス品質特性を評価する枠組みであり、その一つがプロセス能力である。また、組織成熟度の評価にも関係する。対象は、情報技術領域におけるシステムの開発、保守、利用、さらにサービスの設計、移行、提供、改善を含む。

プロセス参照モデルは、ある適用領域のプロセスを、目的と成果、およびプロセス間関係の構造として定義するモデルである。プロセス評価モデルは、一つ以上のプロセス参照モデルに基づき、特定のプロセス品質特性を評価するために適したモデルである。

この種の枠組みは、組織が現在どの程度の能力を持つか、どの能力を強化すべきか、改善が継続的に進んでいるかを考えるために使われる。

\paragraph{表10：枠組みに基づく評価で見る観点}
\begin{longtable}{>{\raggedright\arraybackslash}p{0.28\linewidth}>{\raggedright\arraybackslash}p{0.64\linewidth}}
\toprule
観点 & 説明\\
\midrule
\endfirsthead
\toprule
観点 & 説明\\
\midrule
\endhead
プロセス能力 & 個々のプロセスが目的を達成できる程度である。\\
組織成熟度 & 組織としてプロセスを安定的に運用し改善できる程度である。\\
プロセス参照モデル & プロセスの目的、成果、関係を定義する基準である。\\
プロセス評価モデル & 参照モデルに基づき、特定の品質特性を評価するためのモデルである。\\
\bottomrule
\end{longtable}
\subsection{3.4 アジャイルにおけるプロセス評価と改善}
アジャイル手法では、各反復の終わりにふりかえりを行うことが多い。ふりかえりの目的は、うまくいったこと、うまくいかなかったこと、その理由を分析し、学習と改善のための行動を決めることである。

この実践により、チームは継続的な学習ループの中に入る。名称や範囲は異なるが、このような事後レビューやポストモーテムレビューの考え方は、ソフトウェア工学において以前から存在していた。

重要なのは、ふりかえりを感想会にしないことである。観察事実、測定値、起きた問題、仮説、次の実験を明確にし、次の反復で確かめる必要がある。

\pointbox{ふりかえりの実務ポイント}{「誰が悪いか」ではなく、「どのプロセス条件が望ましくない結果を生んだか」を見る。改善行動は、次の反復で実施できる小さな実験として定義するとよい。}
\section{4 トピックと参考文献の対応}
この表は、SWEBOK が示す「トピック対参考資料の行列」を、学習用に簡略化したものである。詳しく学ぶときは、各トピックに対応する文献や標準を読む。

\paragraph{表11：トピックと主な参考資料}
\begin{longtable}{>{\raggedright\arraybackslash}p{0.28\linewidth}>{\raggedright\arraybackslash}p{0.64\linewidth}}
\toprule
トピック & 主な参考資料\\
\midrule
\endfirsthead
\toprule
トピック & 主な参考資料\\
\midrule
\endhead
1. ソフトウェア工学プロセスの基礎 & ISO/IEC/IEEE 12207、PMBOK Guide、ISO/IEC/IEEE 24765、ISO/IEC 25000、ISO/IEC/IEEE 24748。\\
1.1 導入 & ISO/IEC/IEEE 12207、PMBOK Guide。\\
1.2 プロセス定義 & ISO/IEC/IEEE 12207、ISO/IEC/IEEE 24765、ISO/IEC 25000、ISO/IEC/IEEE 24748、ISO/IEC TR 29110。\\
2. ライフサイクル & Sommerville、Farley、PMI Agile Practice Guide、ISO/IEC/IEEE 32675、ISO/IEC/IEEE 24774。\\
2.1 定義・分類・用語 & ISO/IEC/IEEE 12207、Sommerville、Farley、PMBOK Guide。\\
2.2 ライフサイクルが必要な理由 & Farley、ISO/IEC/IEEE 24774。\\
2.3 プロセスモデルとライフサイクルモデル & Sommerville、PMI Agile Practice Guide、ISO/IEC/IEEE 24765。\\
2.4 開発ライフサイクルモデルの考え方 & Sommerville、Farley、Shore and Warden、PMI Agile Practice Guide、ISO/IEC/IEEE 24765、ISO/IEC/IEEE 32675。\\
2.5 開発ライフサイクルモデルと工学的側面 & Sommerville、Farley、Shore and Warden、PMI Agile Practice Guide、Agile Manifesto、McConnell、Agile Alliance、Eckstein and Buck、RUP、OpenUP。\\
2.6 SLCPの管理 & ISO/IEC/IEEE 24748。\\
2.7 プロセス管理 & ISO/IEC/IEEE 12207、ISO/IEC/IEEE 24765。\\
2.8 ライフサイクル適応 & PMI Software Extension、ISO/IEC/IEEE 24748、ISO/IEC 29110。\\
2.10 プロセス基盤・ツール・手法 & Sommerville、Farley、ISO/IEC/IEEE 24765。\\
2.11 プロセス監視 & ISO/IEC/IEEE 12207、Sommerville、Laporte and April、Farley。\\
3. プロセス評価と改善 & Laporte and April、Shewhart and Deming、ISO/IEC 33001、CMMI、Fenton and Bieman。\\
3.4 アジャイルの評価と改善 & Shore and Warden、Dingsøyr。\\
\bottomrule
\end{longtable}
\section{5 参考文献}
1. ISO/IEC/IEEE 12207:2017, Systems and software engineering — Software life cycle processes.

2. ISO/IEC/IEEE 24765:2017, Systems and Software Engineering — Vocabulary, Edition 2, 2017.

3. I. Sommerville, Software Engineering, 10th ed., 2016.

4. C. Y. Laporte and A. April, Software Quality Assurance, IEEE Computer Society Press, 1st ed., 2018.

5. Project Management Institute, Software Extension to the PMBOK Guide — Fifth Edition, 2013.

6. ISO/IEC 33001:2015, Information technology — Process assessment — Concepts and terminology.

7. ISO/IEC 25000:2014, Systems and software engineering — SQuaRE — Guide to SQuaRE.

8. D. Farley, Modern Software Engineering: Doing What Works to Build Better Software Faster, Addison-Wesley Professional, 2021.

9. J. Shore and S. Warden, The Art of Agile Development, O’Reilly Media, 2nd ed., 2021.

10. Project Management Institute, Agile Practice Guide, Project Management Institute and Agile Alliance, 2017.

11. ISO/IEC/IEEE 32675:2022, Information technology — DevOps — Building reliable and secure systems including application build, package and deployment.

12. ISO/IEC/IEEE 24774:2021, Systems and software engineering — Life cycle management — Specification for process description.

13. Project Management Institute, A Guide to the Project Management Body of Knowledge (PMBOK Guide) — Sixth Edition.

14. ISO/IEC/IEEE 24748-1:2018(E), Systems and software engineering — Life cycle management — Part 1: Guidelines for life cycle management.

15. W. A. Shewhart and W. E. Deming, Statistical Method from the Viewpoint of Quality Control, Dover, 1986.

16. The Agile Manifesto.

17. S. McConnell, More Effective Agile: A Roadmap for Software Leaders, 2019.

18. Agile Alliance, Subway Map to Agile Practices.

19. J. Eckstein and J. Buck, Company-wide Agility with Beyond Budgeting, Open Space \& Sociocracy, 2021.

20. ISO/IEC TR 29110-5-3:2018, Systems and software engineering — Lifecycle profiles for very small entities — Part 5-3.

21. N. Fenton and J. Bieman, Software Metrics, 3rd ed., CRC Press, 2014.

22. CMMI Institute, CMMI V2.0.

23. ISO/IEC/IEEE 24748-3:2020, Guidelines for the application of ISO/IEC/IEEE 12207.

24. D. R. Kiran, Total Quality Management, Elsevier, 2017.

25. J. Rumbaugh, G. Booch, I. Jacobson, The Unified Software Development Process, 1999.

26. P. Kruchten, The Rational Unified Process: An Introduction, 3rd ed., 2004.

27. The Eclipse Foundation.

28. T. Dingsøyr, Postmortem reviews: purpose and approaches in software engineering, Information and Software Technology, 47(5), 2005.

29. ISO/IEC 29110-1-1:2024, Systems and software engineering — Lifecycle profiles for very small entities — Part 1-1: Overview.

30. ISO/IEC/IEEE 31320-1:2012, Information technology — Modeling Languages — Part 1: Syntax and Semantics for IDEF0.

\section{6 画像 TODO 一覧}
本文に配置した画像 TODO を再掲する。実際に図を作る場合は、各 TODO のプロンプトをそのまま画像生成に使うと、A4 資料に合う図を作りやすい。

1. 10章の全体地図を入れる。 画像生成用プロンプト：白背景、モノクロ、A4本文幅に収まる横長の学習用図解を作成する。中央に「ソフトウェア工学プロセス」を置き、周囲に「基礎」「ライフサイクル」「プロセス評価と改善」を大きな三領域として配置する。基礎には「入力・活動・出力」、ライフサイクルには「概念から廃止まで」、評価と改善には「測定・PDCA・GQM」を入れる。日本語ラベル、細線、講義ノート風、余白を広めにする。

2. プロセスの入力・活動・出力の図を入れる。 画像生成用プロンプト：白背景、モノクロ、A4本文幅の横長図を作成する。左に「入力」、中央に「活動」、右に「出力」を置き、矢印で接続する。活動の上に「制約・方針」、下に「人員・ツール・インフラ」を配置し、活動を支えることを示す。出力から次のプロセスの入力へ戻る点線矢印も入れる。日本語ラベル、細線、講義ノート風。

3. 四つのプロセスカテゴリの図を入れる。 画像生成用プロンプト：白背景、モノクロ、A4本文幅の概念図を作成する。中央に「ソフトウェアシステムを生み出すライフサイクル」を置き、周囲に「技術プロセス」「技術管理プロセス」「組織的プロジェクト支援プロセス」「合意プロセス」の四つの箱を配置する。各箱に代表例を二、三個入れる。四つの箱から中央へ矢印を向ける。日本語ラベル、講義ノート風。

4. ライフサイクルモデル比較図を入れる。 画像生成用プロンプト：白背景、モノクロ、A4本文幅の比較表風図を作成する。列に「予測型」「反復型」「進化型」「増分型」「継続的開発」を置く。行に「要求変更への姿勢」「成果物の出し方」「見積りの扱い」「向く状況」を置く。各セルは短い日本語。講義ノート風、細線、印刷で読みやすい。

5. 代表的ライフサイクルモデルの位置づけ図を入れる。 画像生成用プロンプト：白背景、モノクロ、A4本文幅の二軸図を作成する。横軸を「要求変更に閉じている → 変更に開いている」、縦軸を「単一納品に近い → 頻繁なリリースに近い」とする。点として「ウォーターフォール」「Vモデル」「増分」「スパイラル」「迅速プロトタイピング」「統一プロセス」「アジャイル」「DevOps」を配置する。日本語ラベル、講義ノート風。

6. ライフサイクル段階と遷移の図を入れる。 画像生成用プロンプト：白背景、モノクロ、A4本文幅の横長図を作成する。「概念」「開発」「生産」「利用」「支援」「廃止」の箱を並べる。単純な直線矢印だけでなく、利用から支援、支援から開発へ戻る点線矢印も入れ、段階が必ず一方向ではないことを示す。日本語ラベル、講義ノート風。

7. ツール連携としてのソフトウェア工学環境の図を入れる。 画像生成用プロンプト：白背景、モノクロ、A4本文幅のネットワーク図を作成する。中央に「ソフトウェア工学環境」を置き、周囲に「バージョン管理」「テスト」「構成管理」「チケット管理」「ビルド」「要求管理」「測定ダッシュボード」を配置する。各ツール間を細線でつなぎ、「情報を共有する」ことを示す。日本語ラベル、講義ノート風。

8. プロセス監視と製品評価の関係図を入れる。 画像生成用プロンプト：白背景、モノクロ、A4本文幅の循環図を作成する。「プロセス実行」から「成果物・製品」へ矢印を引く。プロセス側に「進捗、欠陥検出率、リードタイム」、製品側に「品質、欠陥、利用者価値」を配置する。両方から「測定・分析」へ矢印を集め、そこから「改善判断」へ、さらにプロセス実行へ戻す。日本語ラベル、講義ノート風。

9. PDCA によるプロセス改善の図を入れる。 画像生成用プロンプト：白背景、モノクロ、A4本文幅の循環図を作成する。「計画」「実行」「確認」「改善」の四つの箱を円形に配置する。中央に「経験的証拠に基づく意思決定」と書く。各段階に短い説明を添える。日本語ラベル、講義ノート風。

10. GQM の階層図を入れる。 画像生成用プロンプト：白背景、モノクロ、A4本文幅の縦長図を作成する。最上段に「目標」、中段に複数の「質問」、下段に複数の「メトリクス」を配置し、上から下へ枝分かれする矢印で接続する。右側に「測定は目標から逆算する」と注記する。日本語ラベル、講義ノート風。

\section{7 章末確認問題}
1. プロセス、活動、タスクの違いを説明せよ。

2. この章でいう「プロセス」が、プログラム実行時の処理ではない理由を説明せよ。

3. ライフサイクルが、単なるプロセス一覧よりも有用である理由を説明せよ。

4. 技術プロセス、技術管理プロセス、組織的プロジェクト支援プロセス、合意プロセスの違いを説明せよ。

5. 予測型、反復型、進化型、増分型、継続的開発を比較せよ。

6. アジャイルは文書不要という意味ではない。なぜそう言えるか。

7. DevOps が必要になった背景を、リリース頻度と組織連携の観点から説明せよ。

8. ライフサイクル適応で、判断と根拠を文書化する理由を説明せよ。

9. GQM の「目標、質問、メトリクス」の順序が重要である理由を説明せよ。

10. アジャイルのふりかえりが、プロセス改善の学習ループである理由を説明せよ。

\pointbox{解答の方向性}{1 は入力・活動・出力と作業粒度の違い、2 は人間の作業活動を扱う点、3 はプロセス間依存と制約を含める点を説明する。5 は要求変更、見積り、成果物の出し方で比較する。9 は測れるものを集めるのではなく、改善目標から測定を設計する点を説明する。}
\end{document}